% TOP_PROBLEM: {"id": 2742, "title": "Kirby Problem 1.83", "category_id": 7, "category": {"id": 7, "name": "topology", "display_name": "Topology", "description": "Properties preserved under continuous deformations.", "slug": "topology", "order_index": 7, "created_at": "2026-07-31T15:26:25.671Z"}, "status": "open"}
% FIRST_POSTED: null
% ENTRY_KIND: statement_only
% Imported from: batch11.tex; UnsolvedMath ID 2742
\documentclass[11pt]{article}
\usepackage[margin=25mm]{geometry}
\usepackage{fontspec}
\setmainfont{Latin Modern Roman}
\setsansfont{Latin Modern Sans}
\usepackage{amsmath,amssymb,mathtools}
\usepackage{unicode-math}
\setmathfont{Latin Modern Math}
\usepackage{xurl,hyperref}
\setcounter{secnumdepth}{0}
\setlength{\parindent}{0pt}
\setlength{\parskip}{5pt}
\setlength{\emergencystretch}{3em}
\newcommand{\sref}[2]{\hyperlink{source:#1:#2}{[S#2]}}
\newcommand{\cataloguescope}[1]{\paragraph{Recorded target.}#1}
\begin{document}
% UnsolvedMath ID 2742; source catalogue code KP-1.83
\setcounter{section}{2741}
\section{Kirby Problem 1.83}\label{2742}
{\small\sffamily UnsolvedMath ID 2742 · Topology}
\subsection{Definitions and mathematical statement}

\paragraph{Shared notation.}$\mathbb N=\{1,2,\ldots\}$, and $\mathbb Z,\mathbb Q,\mathbb R,\mathbb C$ denote the integers, rationals, reals and complex numbers. Any other conventions are specified locally.


An open virtual string is an oriented interval with finitely many pairs of distinct interior points, each pair ordered as an arrow; each marked point belongs to exactly one arrow. It is taken up to orientation-preserving homeomorphism of the interval. These arrows encode self-intersections of a generic oriented path on an oriented surface. Concatenation of intervals gives multiplication. Joining the endpoints gives its closed virtual string $\mu^{\mathrm{cl}}$. Let $\mu'$ reverse both the orientation of the interval and that of every arrow.

A closed string is slice if it is represented by a generic closed curve on the boundary of an oriented 3-manifold that is null-homotopic in that manifold. Use Turaev's cobordism convention: $\mu\sim_c\nu$ precisely when $(\mu\nu')^{\mathrm{cl}}$ is slice. The resulting group is $\mathcal O$, with multiplication induced by concatenation, identity the arrow-free interval, and inverse $[\mu]^{-1}=[\mu']$.
\paragraph{Statement recorded in the supplied source.}
Determine the algebraic structure of $\mathcal O$. In particular, is $\mathcal O$ abelian? Does there exist a nonidentity element $x\in\mathcal O$ and an integer $m\ge2$ with $x^m=1$? \sref{2742}{2} \sref{2742}{3}

\subsection{Short English statement}
Determine the concordance group of open virtual strings, including whether concatenation commutes and whether the group has nontrivial finite-order elements.
\subsection{Sources}
\begin{description}
\item[\hypertarget{source:2742:1}{S1}] ULAM.AI and the UnsolvedMath Contributors, UnsolvedMath: A Curated Collection of Open Mathematics Problems, record 2742 (KP-1.83). CC BY 4.0. \url{https://huggingface.co/datasets/ulamai/UnsolvedMath}.
\item[\hypertarget{source:2742:2}{S2}] R. İnanç Baykur, Robion C. Kirby and Daniel Ruberman, eds., \emph{K3: A New Problem List in Low-Dimensional Topology}, Mathematical Surveys and Monographs 295, American Mathematical Society (2026), Problem 1.83 and its remarks; Chapters 1 and 2 for the stated conventions. \url{https://math.berkeley.edu/sites/default/files/surv-295-ruberman-watermarked-author-pdf.pdf}.
\item[\hypertarget{source:2742:3}{S3}] Vladimir Turaev, \emph{Virtual strings and their cobordisms}, \nolinkurl{arXiv:math/0311185v5} (2004), Sections 5.1, 12.1 and 12.3, and Question 7 in Section 13. \url{https://arxiv.org/pdf/math/0311185v5}.
\end{description}
\label{end:2742}
\end{document}
